\subsection{Approximation of distance functions}\label{subsec:spectral-distance-approximation}
\paraid{p-a700e9644abe}
We first approximate distance profiles in the uniform norm by functions in
finite-dimensional spectral subspaces.
For
$\BasePoint\in\BaseCarrier$,
we call the function
\[
 \MetricSymbol_\BasePoint=\MetricSymbol(\BasePoint,\cdot)
 \in\ContinuousFunctions(\BaseCarrier)
\]
the \emph{distance profile} of
$\BasePoint$.
Its values are the distances from
$\BasePoint$
to all points of
$\BaseCarrier$.
The map
\[
 \BaseCarrier\to\ContinuousFunctions(\BaseCarrier),
 \qquad\BasePoint\longmapsto\MetricSymbol_\BasePoint
\]
is an isometric embedding for the uniform norm,
because
\[
 \norm{\MetricSymbol_\BasePoint-\MetricSymbol_\ComparisonPoint}_\infty
 =\MetricSymbol(\BasePoint,\ComparisonPoint).
\]
This is the distance-function construction of the Kuratowski embedding.
Here the functions
$\MetricSymbol_\BasePoint$
are bounded because
$\MetricSpace{\BaseCarrier}{\MetricSymbol}$
is compact.
The following theorem approximates all distance profiles in one finite
spectral subspace.

\begin{theorem}\label{lem:spectral-density}
\paraid{p-c4a771b9c692}
For every nonempty compact metric space
$\MetricSpace{\BaseCarrier }{\MetricSymbol }$
and every
$\SmallError >0$,
there is
$\SpectralCutoff >0$
such that
$\SpectralCutoff ,-\SpectralCutoff \notin\spec(\DistanceOperator{\BaseCarrier}{\MetricSymbol})$
and
\begin{equation}\label{eq:spectral-uniform-approximation}
 \sup_{\BasePoint \in \BaseCarrier }\inf_{\Vector \in\SpectralSubspace{\BaseCarrier }{\SpectralCutoff }}
          \norm{\MetricSymbol _\BasePoint -\Vector }_\infty<\SmallError .
\end{equation}
If $\Card(\BaseCarrier)\geq2$,
then we may also require
$\dim\SpectralSubspace{\BaseCarrier }{\SpectralCutoff }\geq1$.
\end{theorem}

\begin{proof}
\paraid{p-53612876bfd1}
Let
$\UniformSpectralSpan $
be the closed linear span in
$\ContinuousFunctions(\BaseCarrier )$
of the continuous representatives of all nonzero eigenfunctions.
For
$\HilbertFunction \in \HilbertDomain _\BaseCarrier $,
write
$\HilbertFunction =\HilbertFunction _0+\HilbertFunction _\perp$
with
$\HilbertFunction _0\in\ker\DistanceOperator{\BaseCarrier}{\MetricSymbol}$
and
$\HilbertFunction _\perp\perp\ker\DistanceOperator{\BaseCarrier}{\MetricSymbol}$.
Choose a decreasing sequence of positive real numbers
$\{\SpectralCutoff_m\}_{m\in\NonnegativeIntegers}$
converging to
$0$
such that, for every $m\in\NonnegativeIntegers$,
$\SpectralCutoff_m,-\SpectralCutoff_m\notin\spec(\DistanceOperator{\BaseCarrier}{\MetricSymbol})$.
Such a sequence exists because the nonzero spectrum is countable.
By the Riesz--Schauder and self-adjoint spectral theorem (\Cref{prop:spectral-basics}),
we have
$\norm{\SpectralProjection{\BaseCarrier }{\SpectralCutoff_m}\HilbertFunction -\HilbertFunction _\perp}_2\to0$.
Since
$\HilbertFunction_0\in\ker\DistanceOperator{\BaseCarrier}{\MetricSymbol}$,
the continuous function
$\DistanceOperator{\BaseCarrier}{\MetricSymbol}\HilbertFunction_0$
vanishes almost everywhere.
By full support of
$\SelectedMeasure{\BaseCarrier}{\MetricSymbol}$,
for every
$\BasePoint\in\BaseCarrier$
we have
\begin{equation}\label{eq:continuous-kernel-vanishing}
 (\DistanceOperator{\BaseCarrier}{\MetricSymbol}\HilbertFunction_0)(\BasePoint)=0.
\end{equation}
Using \eqref{eq:operator-uniform-bound-1},
\[
 \norm{\DistanceOperator{\BaseCarrier}{\MetricSymbol}\SpectralProjection{\BaseCarrier }{\SpectralCutoff_m}\HilbertFunction
                 -\DistanceOperator{\BaseCarrier}{\MetricSymbol}\HilbertFunction }_\infty
 \leq\diam\MetricSpace{\BaseCarrier }{\MetricSymbol }
       \norm{\SpectralProjection{\BaseCarrier }{\SpectralCutoff_m}\HilbertFunction -\HilbertFunction _\perp}_2\to0.
\]
For every
$\SpectralCutoff>0$,
we have
\[
 \DistanceOperator{\BaseCarrier}{\MetricSymbol}
 \SpectralProjection{\BaseCarrier}{\SpectralCutoff}\HilbertFunction
 \in\SpectralSubspace{\BaseCarrier}{\SpectralCutoff}
 \subset\UniformSpectralSpan.
\]
Since
$\UniformSpectralSpan$
is closed in the uniform norm,
it follows that
$\DistanceOperator{\BaseCarrier}{\MetricSymbol}\HilbertFunction \in \UniformSpectralSpan $.

\paraid{p-4c9880842102}
Fix
$\BasePoint \in \BaseCarrier $
and
$\Radius >0$.
Use the normalized function
$\NormalizedBump{\BaseCarrier}{\MetricSymbol}{\BasePoint}{\Radius}{\SelectedMeasure{\BaseCarrier}{\MetricSymbol}}$
from \eqref{eq:normalized-bump-definition}.
By \eqref{eq:normalized-bump-properties-1},
\eqref{eq:normalized-bump-properties-2} and \eqref{eq:normalized-bump-properties-3},
the normalized function is nonnegative,
has integral one,
and is supported in
$\MetricBall(\BasePoint,\Radius;\MetricSymbol)$.
For every
$\ComparisonPoint \in \BaseCarrier $,
\begin{equation}\label{eq:bump-distance-approximation}
 \begin{aligned}
 &|\DistanceOperator{\BaseCarrier}{\MetricSymbol}\NormalizedBump{\BaseCarrier}{\MetricSymbol}{\BasePoint}{\Radius}{\SelectedMeasure{\BaseCarrier}{\MetricSymbol}}(\ComparisonPoint )-\MetricSymbol _\BasePoint (\ComparisonPoint )|\\
 &=\left|\int_\BaseCarrier (\MetricSymbol (\ComparisonPoint ,\IntegrationPoint )-\MetricSymbol (\ComparisonPoint ,\BasePoint ))\NormalizedBump{\BaseCarrier}{\MetricSymbol}{\BasePoint}{\Radius}{\SelectedMeasure{\BaseCarrier}{\MetricSymbol}}(\IntegrationPoint )
                                      \,\IntegrationDifferential \SelectedMeasure{\BaseCarrier}{\MetricSymbol}(\IntegrationPoint )\right|\\
 &\leq\int_\BaseCarrier  \MetricSymbol (\BasePoint ,\IntegrationPoint )\NormalizedBump{\BaseCarrier}{\MetricSymbol}{\BasePoint}{\Radius}{\SelectedMeasure{\BaseCarrier}{\MetricSymbol}}(\IntegrationPoint )\,\IntegrationDifferential \SelectedMeasure{\BaseCarrier}{\MetricSymbol}(\IntegrationPoint )
 \leq \Radius .
 \end{aligned}
\end{equation}
Since \eqref{eq:bump-distance-approximation} holds for every
$\Radius>0$
and
$\UniformSpectralSpan$
is closed in the uniform norm,
we obtain
$\MetricSymbol_\BasePoint\in\UniformSpectralSpan$.

\paraid{p-0c97576ebc0b}
Recall that
\[
 \norm{\MetricSymbol _\BasePoint -\MetricSymbol _\ComparisonPoint }_\infty=\MetricSymbol (\BasePoint ,\ComparisonPoint ),
\]
as established before the theorem.
Choose finitely many points
$\BasePoint_0,\ldots,\BasePoint_{\NetSize-1}$
whose
$\SmallError /3$-balls
cover
$\BaseCarrier $.
For each
$\NetIndex $,
choose a finite sum
$\Vector _\NetIndex $
of nonzero eigenfunctions with
$\norm{\MetricSymbol _{\BasePoint _\NetIndex }-\Vector _\NetIndex }_\infty<\SmallError /3$.
Let
$\ApproximationEigenvalues$
be the finite set of eigenvalues occurring in these sums.
Choose
$\SpectralCutoff>0$
such that
\[
 \SpectralCutoff,-\SpectralCutoff\notin\spec(\DistanceOperator{\BaseCarrier}{\MetricSymbol}),
\]
and such that for every
$\Eigenvalue\in\ApproximationEigenvalues$
we have
\[
 \SpectralCutoff<|\Eigenvalue|.
\]
Note that even if
$\ApproximationEigenvalues=\emptyset$,
it suffices to choose a positive real number
$\SpectralCutoff$
satisfying
$\SpectralCutoff,-\SpectralCutoff\notin\spec(\DistanceOperator{\BaseCarrier}{\MetricSymbol})$.
Such a real number exists because the nonzero spectrum is countable.
Then every
$\Vector _\NetIndex $
belongs to
$\SpectralSubspace{\BaseCarrier }{\SpectralCutoff }$.
For
$\MetricSymbol (\BasePoint ,\BasePoint _\NetIndex )<\SmallError /3$,
\[
 \inf_{\Vector \in\SpectralSubspace{\BaseCarrier }{\SpectralCutoff }}\norm{\MetricSymbol _\BasePoint -\Vector }_\infty
 \leq\norm{\MetricSymbol _\BasePoint -\MetricSymbol _{\BasePoint _\NetIndex }}_\infty+\norm{\MetricSymbol _{\BasePoint _\NetIndex }-\Vector _\NetIndex }_\infty
 <2\SmallError /3.
\]
This proves \eqref{eq:spectral-uniform-approximation}.

\paraid{p-f6773180e565}
If
$\BaseCarrier $
is not a singleton,
choose
$\BasePoint ,\IntegrationPoint \in \BaseCarrier $
with
$\MetricSymbol (\BasePoint ,\IntegrationPoint )>0$.
The integral
$\DistanceOperator{\BaseCarrier}{\MetricSymbol}1(\BasePoint )$
is positive because the distance from
$\BasePoint $
is positive on a neighborhood of
$\IntegrationPoint $
of positive measure.
Thus
$\DistanceOperator{\BaseCarrier}{\MetricSymbol}\ne0$.
By the Riesz--Schauder and self-adjoint spectral theorem (\Cref{prop:spectral-basics}),
it has a nonzero eigenvalue.
If necessary,
choose a smaller positive
$\SpectralCutoff$
below the absolute value of this eigenvalue and still satisfying
$\pm\SpectralCutoff\notin\spec(\DistanceOperator{\BaseCarrier}{\MetricSymbol})$.
This retains the approximation in \eqref{eq:spectral-uniform-approximation} and ensures
$\dim\SpectralSubspace{\BaseCarrier}{\SpectralCutoff}\geq1$.
\end{proof}

\paraid{p-9e170f5d7db0}
The uniform distance between the distance profiles is the original distance,
but a best approximation in a finite-dimensional subspace need not be
unique for the uniform norm.
We minimize a finite
$\LebesgueSpace^\IntegrabilityExponent $-norm
instead.
For
$1<\IntegrabilityExponent <\infty$,
the best approximation exists and is unique by \Cref{thm:best-approximation}.
\Cref{lem:lp-distance} shows that the
$\LebesgueSpace^\IntegrabilityExponent$-distances
between distance profiles
approximate the original distances uniformly when
$\IntegrabilityExponent$
is sufficiently large.
We then combine \eqref{eq:lp-distance-error} with finite spectral approximation in
\eqref{eq:finite-spectral-error}.

\paraid{p-c8d27acdc339}
To compare the best approximations,
we use the following continuity statement for their coefficients.

\begin{lemma}\label{lem:minimizer-continuity}
\paraid{p-df64ce5d76d4}
Let
$\ModelDimension\geq1$,
and let continuous functions
$\ApproximationObjective_\SequenceIndex,\ApproximationObjective\colon
 \RealNumbers^\ModelDimension\to\RealNumbers$
converge uniformly on every compact subset.
Assume that each
$\CoefficientVector_\SequenceIndex$
minimizes
$\ApproximationObjective_\SequenceIndex$,
that these vectors form a bounded sequence,
and that
$\ApproximationObjective$
has a unique minimizer
$\CoefficientVector$.
Then
$\CoefficientVector_\SequenceIndex\to\CoefficientVector$.
\end{lemma}

\begin{proof}
\paraid{p-586499da67ea}
By boundedness of $\CoefficientVector_j$,
choose
$\CoefficientBound>0$
such that
\[
 \{\CoefficientVector_\SequenceIndex\mid\SequenceIndex\in\NonnegativeIntegers\}
 \subset\{\CompetitorCoefficients\in\RealNumbers^\ModelDimension
       \mid\norm{\CompetitorCoefficients}_2\leq\CoefficientBound\}.
\]
This closed ball is compact,
so every subsequence has a further convergent subsequence.
Let
$\ComparisonCoefficients$
denote the limit of such a subsequence
$\{\CoefficientVector_{n_k}\}_{k\in\NonnegativeIntegers}$.
The limit satisfies
$\norm{\ComparisonCoefficients}_2\leq\CoefficientBound$,
and the assumed uniform convergence implies
\begin{equation}\label{eq:minimizer-uniform-objective-error}
 \sup_{\substack{\CompetitorCoefficients\in\RealNumbers^\ModelDimension\\\norm{\CompetitorCoefficients}_2\leq\CoefficientBound}}
 |\ApproximationObjective_{n_k}(\CompetitorCoefficients)
       -\ApproximationObjective(\CompetitorCoefficients)|\longrightarrow0.
\end{equation}
Since $\ApproximationObjective$ is continuous and
$\CoefficientVector_{n_k}\to\ComparisonCoefficients$,
the difference between its values at these two coefficient vectors tends to $0$.
The objective error
$|\ApproximationObjective_{n_k}(\CoefficientVector_{n_k})
 -\ApproximationObjective(\CoefficientVector_{n_k})|$
tends to $0$ by \eqref{eq:minimizer-uniform-objective-error}.
Consequently,
\begin{equation}\label{eq:minimizer-objective-limit}
 \begin{aligned}
 &|\ApproximationObjective_{n_k}(\CoefficientVector_{n_k})
        -\ApproximationObjective(\ComparisonCoefficients)|
 \leq\\
 &|\ApproximationObjective_{n_k}(\CoefficientVector_{n_k})
        -\ApproximationObjective(\CoefficientVector_{n_k})|+|\ApproximationObjective(\CoefficientVector_{n_k})
        -\ApproximationObjective(\ComparisonCoefficients)|\longrightarrow0.
 \end{aligned}
\end{equation}
For a fixed competitor
$\CompetitorCoefficients\in\RealNumbers^\ModelDimension$,
minimality at each index and convergence at that competitor now imply
\begin{equation}\label{eq:limit-minimizer-comparison}
 \ApproximationObjective(\ComparisonCoefficients)
 =\lim_{k\to\infty}\ApproximationObjective_{n_k}(\CoefficientVector_{n_k})
 \leq\lim_{k\to\infty}\ApproximationObjective_{n_k}(\CompetitorCoefficients)
 =\ApproximationObjective(\CompetitorCoefficients).
\end{equation}
Thus
$\ComparisonCoefficients=\CoefficientVector$
by uniqueness.
Relative compactness and uniqueness of the subsequential limit imply
$\CoefficientVector_\SequenceIndex\to\CoefficientVector$.
\end{proof}

\begin{lemma}\label{lem:lp-distance}
\paraid{p-c7556a5ad9a5}
Let
$\MetricSpace{\BaseCarrier }{\MetricSymbol }$
be a nonempty compact metric space and set
$\DiameterBound =\diam\MetricSpace{\BaseCarrier }{\MetricSymbol }$.
Then for every
$\Radius >0$,
the number
\[
 \BallMassBound _\Radius =\inf_{\BasePoint \in \BaseCarrier }\SelectedMeasure{\BaseCarrier}{\MetricSymbol}(\MetricBall(\BasePoint ,\Radius;\MetricSymbol))
\]
is positive.
Moreover,
for every
$\Radius>0$,
every
$2\leq \IntegrabilityExponent <\infty$,
and every pair
$\BasePoint ,\ComparisonPoint \in \BaseCarrier $,
\begin{equation}\label{eq:lp-distance-error}
 0\leq \MetricSymbol (\BasePoint ,\ComparisonPoint )-\norm{\MetricSymbol _\BasePoint -\MetricSymbol _\ComparisonPoint }_{\LebesgueSpace^\IntegrabilityExponent (\SelectedMeasure{\BaseCarrier}{\MetricSymbol})}
 \leq \DiameterBound (1-\BallMassBound _\Radius ^{1/\IntegrabilityExponent })+2\Radius .
\end{equation}
\end{lemma}

\begin{proof}
\paraid{p-ce0329a6de36}
In this proof, $\norm{\cdot}_\IntegrabilityExponent$ denotes the norm on
$\LebesgueSpace^\IntegrabilityExponent(\BaseCarrier,\SelectedMeasure{\BaseCarrier}{\MetricSymbol})$.
Choose a finite
$\Radius /4$-net
$\{\IntegrationPoint_0,\ldots,\IntegrationPoint_{\NetSize-1}\}$.
For each
$\BasePoint \in \BaseCarrier $,
some
$\NetIndex $
satisfies
$\MetricSymbol (\BasePoint ,\IntegrationPoint _\NetIndex )<\Radius /4$,
and then
$\MetricBall(\IntegrationPoint _\NetIndex ,\Radius /4;\MetricSymbol)\subset\MetricBall(\BasePoint ,\Radius;\MetricSymbol)$.
Therefore,
\[
 \BallMassBound _\Radius \geq\min_{0\leq \NetIndex<\NetSize }
          \SelectedMeasure{\BaseCarrier}{\MetricSymbol}(\MetricBall(\IntegrationPoint _\NetIndex ,\Radius /4;\MetricSymbol))>0.
\]
For all
$\IntegrationPoint \in \BaseCarrier $,
$|\MetricSymbol _\BasePoint (\IntegrationPoint )-\MetricSymbol _\ComparisonPoint (\IntegrationPoint )|\leq \MetricSymbol (\BasePoint ,\ComparisonPoint )$.
For
$\IntegrationPoint \in\MetricBall(\BasePoint ,\Radius;\MetricSymbol)$,
\[
 \MetricSymbol _\ComparisonPoint (\IntegrationPoint )-\MetricSymbol _\BasePoint (\IntegrationPoint )\geq \MetricSymbol (\BasePoint ,\ComparisonPoint )-2\MetricSymbol (\BasePoint ,\IntegrationPoint )>\MetricSymbol (\BasePoint ,\ComparisonPoint )-2\Radius .
\]
Consequently,
\begin{equation}\label{eq:lp-distance-two-sided-bound}
 \BallMassBound _\Radius ^{1/\IntegrabilityExponent }\max\{\MetricSymbol (\BasePoint ,\ComparisonPoint )-2\Radius ,0\}
 \leq\norm{\MetricSymbol _\BasePoint -\MetricSymbol _\ComparisonPoint }_\IntegrabilityExponent\leq \MetricSymbol (\BasePoint ,\ComparisonPoint ).
\end{equation}
Since
$\SelectedMeasure{\BaseCarrier}{\MetricSymbol}$
is a probability and the ball masses are positive,
$0<\BallMassBound_\Radius\leq1$.
Consequently,
$0<\BallMassBound_\Radius^{1/\IntegrabilityExponent}\leq1$.
From the lower bound in \eqref{eq:lp-distance-two-sided-bound} and
\[
\max\{\MetricSymbol(\BasePoint,\ComparisonPoint)-2\Radius,0\}
 \geq\MetricSymbol(\BasePoint,\ComparisonPoint)-2\Radius
 \]
it follows that
\[
 \begin{aligned}
 \MetricSymbol(\BasePoint,\ComparisonPoint)
   -\norm{\MetricSymbol_\BasePoint-\MetricSymbol_\ComparisonPoint}_\IntegrabilityExponent
 &\leq \MetricSymbol(\BasePoint,\ComparisonPoint)
 -\BallMassBound_\Radius^{1/\IntegrabilityExponent}
 \max\{\MetricSymbol(\BasePoint,\ComparisonPoint)-2\Radius,0\}\\
 &\leq \MetricSymbol(\BasePoint,\ComparisonPoint)
 -\BallMassBound_\Radius^{1/\IntegrabilityExponent}
 \bigl(\MetricSymbol(\BasePoint,\ComparisonPoint)-2\Radius\bigr)\\
 &= (1-\BallMassBound_\Radius^{1/\IntegrabilityExponent})
      \MetricSymbol(\BasePoint,\ComparisonPoint)
       +2\Radius\BallMassBound_\Radius^{1/\IntegrabilityExponent}\\
 &\leq
 (1-\BallMassBound_\Radius^{1/\IntegrabilityExponent})
      \DiameterBound+2\Radius.
 \end{aligned}
\]
The last inequality uses
$\MetricSymbol(\BasePoint,\ComparisonPoint)\leq\DiameterBound$,
$1-\BallMassBound_\Radius^{1/\IntegrabilityExponent}\geq0$,
and
$\BallMassBound_\Radius^{1/\IntegrabilityExponent}\leq1$.
This proves \eqref{eq:lp-distance-error}.
\end{proof}

\paraid{p-533de80974c4}
Fix
$\SpectralSubspace{\BaseCarrier }{\SpectralCutoff }$
and
$2\leq \IntegrabilityExponent <\infty$.
Write $\norm{\cdot}_\IntegrabilityExponent$ for the norm on
$\LebesgueSpace^\IntegrabilityExponent(\BaseCarrier,\SelectedMeasure{\BaseCarrier}{\MetricSymbol})$.
For each
$\BasePoint \in \BaseCarrier $,
let
$\Vector _\BasePoint $
be the unique minimizer of
$\Vector \mapsto\norm{\MetricSymbol _\BasePoint -\Vector }_\IntegrabilityExponent ^\IntegrabilityExponent $
over
$\SpectralSubspace{\BaseCarrier }{\SpectralCutoff }$.
The spectral subspace
$\SpectralSubspace{\BaseCarrier}{\SpectralCutoff}$
is finite-dimensional,
and hence closed and convex in
$\LebesgueSpace^\IntegrabilityExponent(\BaseCarrier,\SelectedMeasure{\BaseCarrier}{\MetricSymbol})$.
By \Cref{thm:best-approximation},
this minimizer exists and is unique.
Since
$\SelectedMeasure{\BaseCarrier}{\MetricSymbol}$
has full support,
the natural map
$\ContinuousFunctions(\BaseCarrier)\to
 \LebesgueSpace^\IntegrabilityExponent(\BaseCarrier,\SelectedMeasure{\BaseCarrier}{\MetricSymbol})$
is injective.

\paraid{p-f10c0269a0dd}
We call the function
$\Vector_\BasePoint\in\SpectralSubspace{\BaseCarrier}{\SpectralCutoff}$
the \emph{finite-dimensional approximation of the distance profile}.
The map
$\BaseCarrier\to\SpectralSubspace{\BaseCarrier}{\SpectralCutoff}$,
$\BasePoint\mapsto\Vector_\BasePoint$,
need not be injective.
We use these approximating functions to approximate distances with a controlled error.

\paraid{p-774eca26292a}
If \eqref{eq:spectral-uniform-approximation} holds,
then
$\norm{\MetricSymbol _\BasePoint -\Vector _\BasePoint }_\IntegrabilityExponent <\SmallError $.
Define
$\Pseudometric_\BaseCarrier\colon\BaseCarrier\times\BaseCarrier\to[0,\infty)$
by
$\Pseudometric_\BaseCarrier(\BasePoint,\ComparisonPoint)
 =\norm{\Vector_\BasePoint-\Vector_\ComparisonPoint}_\IntegrabilityExponent$.
The norm triangle inequality implies
\[
 \left|\Pseudometric _\BaseCarrier (\BasePoint ,\ComparisonPoint )-\norm{\MetricSymbol _\BasePoint -\MetricSymbol _\ComparisonPoint }_\IntegrabilityExponent \right|
 \leq\norm{\Vector _\BasePoint -\MetricSymbol _\BasePoint }_\IntegrabilityExponent +\norm{\Vector _\ComparisonPoint -\MetricSymbol _\ComparisonPoint }_\IntegrabilityExponent <2\SmallError .
\]
Together with \Cref{lem:lp-distance},
this yields
\begin{equation}\label{eq:finite-spectral-error}
 \norm{\Pseudometric _\BaseCarrier -\MetricSymbol }_\infty
 \leq2\SmallError +\DiameterBound (1-\BallMassBound _\Radius ^{1/\IntegrabilityExponent })+2\Radius .
\end{equation}
For a prescribed error tolerance,
choose
$\Radius $
small,
then choose a sufficiently large finite
$\IntegrabilityExponent $,
and finally choose
$\SmallError $
small and a corresponding spectral subspace.
